\chapter{环境相关疾病}
\setcounter{chapter}{8}
\chapterintro{
掌握生物地球化学性疾病的概念及流行特征；掌握碘缺乏病、地方性氟中毒、地方性砷中毒的病因、发病机制、临床表现及预防措施；掌握克山病与大骨节病；了解水俣病和宣威肺癌的环境流行病学特征。
}

\section{概述}

\begin{defbox}
\textbf{生物地球化学性疾病（Biogeochemical Disease）}：由于地壳表面化学元素分布的不均匀性，使某些地区的水和/或土壤中某些元素过多或过少，当地居民通过饮水、食物等途径摄入这些元素过多或过少，而引起的特异性疾病。

\textbf{地方病（Endemic Disease）}：由于自然因素和社会因素的影响，在某一地区的人群中发生，不需自外地输入，并呈地方性流行特点的某种疾病。
\end{defbox}

\textbf{流行特征：}①\imp{明显的地区性分布}；②\imp{与环境中元素水平相关}。\textbf{影响因素：}营养条件、生活习惯、多种元素的联合作用。

\textbf{我国主要防治的生物地球化学性疾病（5种）：}碘缺乏病、地方性氟中毒、地方性砷中毒、克山病、大骨节病。

\textbf{综合控制措施：}①限制摄入（改换水源除氟/砷、改良炉灶、改变生活方式）；②适量补充（食盐加碘——最安全有效、口服亚硒酸钠补硒）；③组织措施（建立防治机构、监测病情、动员群众参与）。

\section{碘与健康——碘缺乏病}

\begin{defbox}
\textbf{碘缺乏病（IDD）}：从胚胎发育至成人期由于碘摄入量不足而引起的一系列病症。包括地方性甲状腺肿、地方性克汀病（呆小症）、地方性亚临床克汀病、碘缺乏导致的流产、早产、死产、先天畸形等。

\textbf{地方性克汀病（呆小症）}：由于外环境较严重缺碘引起的以脑发育障碍和体格发育落后为主要特征的地方病。主要表现为较严重的智力障碍、聋哑、神经运动功能障碍、体格发育落后等——\imp{"呆、小、聋、哑、瘫"}。

\textbf{剂量参考：}最低生理需要量75 $\mu$g/d，供给量150 $\mu$g/d，孕妇250 $\mu$g/d；摄入碘低于40 $\mu$g/d、水中含碘量低于10 $\mu$g/L时，碘缺乏病流行。
\end{defbox}

\subsection{流行特征}

\begin{itemize}[leftmargin=*]
    \item 地区分布：\imp{山区 $>$ 丘陵 $>$ 平原 $>$ 沿海}；内陆 $>$ 沿海
    \item 人群分布：任何年龄均可发病，女性早于男性
    \item 时间趋势：补碘干预后可迅速改变流行状况
\end{itemize}

\subsection{发病机制}

碘缺乏 → 甲状腺激素合成不足 → TSH反馈性升高 → 甲状腺代偿性增生 → \imp{地方性甲状腺肿}。重度缺碘影响胎儿脑发育 → \imp{地方性克汀病}（呆小症）。

\subsection{预防措施}

\imp{碘盐}（最经济有效）；碘油（重点人群）；富碘食物（海带、紫菜等）。

\section{氟与健康——地方性氟中毒}

\begin{defbox}
\textbf{地方性氟中毒（Endemic Fluorosis）}：由于一定地区的环境中氟元素过多，导致生活在 该环境中的居民经饮水、食物和空气等途径长期摄入过量氟，引起以\keyword{氟骨症和氟斑牙}为主要特征的慢性全身性疾病。人体摄入总氟量超过4 mg/d即可引起慢性氟中毒。
\end{defbox}

\subsection{发病机制详解}

\begin{enumerate}[leftmargin=*]
    \item \imp{对骨组织的影响}：羟基磷灰石→氟磷灰石→难溶性氟化钙（CaF$_2$）。氟对成骨细胞的作用呈双向性：长期小剂量氟促进骨形成，大剂量诱导细胞凋亡引起骨质疏松。双向作用使氟中毒时出现骨质硬化、骨质疏松或两者同时并存
    \item \imp{对钙磷代谢的影响}：过量氟消耗大量钙→血钙降低→刺激甲状旁腺分泌激素增多→抑制肾小管对磷的重吸收→磷排出增多→磷代谢紊乱。血钙减少和甲状旁腺激素增加又刺激钙从骨组织中不断释放入血，造成骨质脱钙或溶骨，临床上可表现为骨质疏松、骨软化甚至骨骼变形
    \item \imp{抑制酶活性}：氟可与某些酶结构中的金属离子形成复合物，或与带正电的赖氨酸和精氨酸基团、磷蛋白等结合，改变酶结构、抑制酶活性
    \item \imp{对牙齿的影响——氟斑牙}：氟化钙沉积于发育中的牙组织，导致色素沉着、易碎缺损
    \item \imp{非骨相损害}：以中枢神经系统损害最为常见；损伤神经受体；使骨骼肌纤维萎缩；损伤肾小管；诱导肝细胞凋亡；使血管壁钙化硬化；使C细胞功能紊乱；破坏睾丸细胞结构导致生殖功能下降
\end{enumerate}

\subsection{氟的双重健康效应}

\begin{itemize}[leftmargin=*]
    \item \imp{适量}：预防龋齿，促进骨骼钙化
    \item \imp{过量}：氟斑牙、氟骨症
    \item \imp{缺乏}：龋齿发病率增加
\end{itemize}

\subsection{病区类型与预防}

\begin{table}[htbp]
\centering
\caption{地方性氟中毒病区类型及预防措施}
\begin{tabularx}{\textwidth}{lX}
\hline
\textbf{病区类型} & \textbf{预防措施} \\
\hline
\imp{饮水型}（分布最广） & 改换水源（改用低氟水源）；饮水除氟 \\
\imp{燃煤污染型} & 改良炉灶；减少食物氟污染；不用高氟劣质煤 \\
\imp{饮砖茶型} & 研制低氟砖茶；降低砖茶氟含量 \\
\hline
\end{tabularx}
\end{table}

\subsection{人群分布}

年龄：氟斑牙——生长发育中的恒牙（7--8岁前）；氟骨症——成年人。性别：女性以骨软化为主，男性以骨硬化为主。居住时间越长病情越重。

\section{砷与健康——地方性砷中毒}

\begin{defbox}
\textbf{地方性砷中毒（Endemic Arsenism）}：长期从饮用水、室内煤烟、食物等环境介质中摄入过量砷引起以\keyword{皮肤损害}（色素沉着、脱失、角化）为主要特征，伴有多系统、多脏器受损的慢性全身性疾病。砷为IARC确认的\imp{1类致癌物}。

\textbf{砷的代谢：}
\begin{itemize}[leftmargin=*]
    \item 吸收：呼吸道（三价砷，以颗粒物为载体沉积肺组织）；消化道（吸收容易度：五价砷$>$三价砷，无机砷$>$有机砷）；皮肤黏膜（贮存于角蛋白中）
    \item 分布蓄积：三价砷极易与巯基结合→分布于肝、肾、脑等实质性脏器；五价砷取代骨组织中磷灰石的磷酸盐→蓄积于骨组织；三价砷易与角蛋白结合→蓄积于皮肤、指（趾）甲、毛发。毛发砷含量为早期、敏感的内暴露生物标志
    \item 排泄：肾脏是主要排泄器官，尿砷测定可灵敏反映机体砷内暴露水平
\end{itemize}

\textbf{发病机制：}①抑制酶活性——三价砷与酶蛋白双巯基或羧基结合形成稳定络合物，五价砷与酶结合物不稳定可自行水解恢复活性；砷对酪氨酸酶呈双相反应——小剂量产生大量黑色素，大量蓄积致皮肤色素脱失；②导致细胞凋亡；③致癌——直接导致DNA损伤、影响损伤修复、基因表达异常、DNA甲基化反应、产生活性氧引起氧化应激。

\textbf{病区类型：}饮水型（台湾"乌脚病"、新疆、内蒙古）和燃煤污染型（我国特有，分布于贵州、四川、陕西，陕西为"氟砷联合污染型"）。
\end{defbox}

\section{硒与健康}

\subsection{克山病}

\begin{defbox}
\textbf{克山病（Keshan Disease）}：与环境低硒有关的地方性心肌病。1935年首先发现于黑龙江省克山县。
\end{defbox}

\textbf{病因：}环境硒缺乏 + 柯萨奇病毒感染等多因素。\textbf{病理：}心肌变性、坏死、纤维化。\textbf{预防：}口服\imp{亚硒酸钠}补硒，改善膳食营养。

\subsection{大骨节病}

\textbf{大骨节病（Kaschin-Beck Disease）}：与环境低硒有关的慢性骨关节病。以骺板软骨和关节软骨变性坏死为特征，导致关节增粗、短指畸形、身材矮小。

\section{环境污染性疾病}

\begin{defbox}
\textbf{环境污染性致病因素}：能污染环境，使环境质量恶化，而直接或间接使人患病的环境污染因素。

\textbf{环境污染性疾病}：由环境污染因素在暴露人群中引发的疾病。

\textbf{公害病}：严重的地区性环境污染性疾病，由政府认定，须经严格鉴定和国家法律的正式认可，有\imp{医学和法律的双重含义}，病人享受国家补贴。
\end{defbox}

\textbf{环境污染性疾病的特点：}
\begin{enumerate}[leftmargin=*]
    \item 环境污染区域内的人群不分年龄、性别都可能发病
    \item 发病者均出现与暴露污染物相关的相同症状和体征
    \item 除急性危害外，大多具有低浓度长期暴露、陆续发病的特点
    \item 往往缺乏健康危害的早期诊断指标
    \item 预防的关键在于消除环境污染性致病因素、加强对易感人群和亚临床阶段人群的保护
\end{enumerate}

\subsection{慢性甲基汞中毒（水俣病）}

\begin{keybox}
\textbf{水俣病——慢性甲基汞中毒：}
\begin{itemize}[leftmargin=*]
    \item 污染来源：含汞工业废水→无机汞在底泥中经微生物作用甲基化→甲基汞经食物链生物放大（鱼贝类富集百万倍以上）→人食鱼中毒
    \item 发病机制：甲基汞主要经消化道摄入，具有脂溶性、原形蓄积和高神经毒特性。甲基汞与胃酸作用生成氯化甲基汞，吸收入血后与血红蛋白中的巯基结合，随血液分布到脑、肝、肾脏和其他组织。甲基汞能通过血脑屏障进入脑细胞，还能透过胎盘进入胎儿脑中。脑细胞富含类脂质，甲基汞对此有很高亲和力，易蓄积在脑细胞内，造成\imp{进行性不可逆性}的病理损害
    \item 临床表现：感觉障碍（口周和肢端麻木）、共济失调、视野缩小、听力障碍。先天性水俣病：脑性瘫痪——甲基汞通过胎盘导致胎儿中枢神经系统发育障碍
\end{itemize}
\end{keybox}

\subsection{宣威肺癌与军团菌病}

\subsection{军团病}

\begin{keybox}
\textbf{军团病（Legionnaires Disease, LD）}：由\imp{嗜肺军团菌（Legionella Pneumophila, LP）}引起的一种以肺炎为主的全身性疾病，以肺部感染伴全身多系统损伤为主要表现，也可表现为无肺炎的急性自限性流感样疾病。

\textbf{流行病学特征：}
\begin{itemize}[leftmargin=*]
    \item \textbf{传染源}：人工水环境，包括冷热水管道系统、集中空调冷却塔循环水、淋浴设施、游泳池、喷泉、空气加湿器
    \item \textbf{传播途径}：吸入被军团菌污染的气溶胶
    \item \textbf{时间分布}：多发于夏秋季
    \item \textbf{人群分布}：中老年男性和烟酒嗜好者多见
\end{itemize}

\textbf{诊断依据：}
\begin{enumerate}[leftmargin=*]
    \item 流行病学资料：夏秋发病，环境中有建筑施工、空调系统、淋浴喷头设施
    \item 临床资料：临床表现+胸片+非特异性实验室检查（无特异性、多样性）；特异性检查——呼吸道分泌物和肺组织标本培养、血清/呼吸道分泌物抗体荧光抗体法检测、痰/尿液抗原ELISA检测
    \item 环境检测：对室内空气、冷凝塔水等环境样品进行细菌培养
\end{enumerate}
\end{keybox}

\begin{tipbox}
\textbf{宣威女性肺癌高发：}室内燃煤产生PAHs等致癌物→女性暴露更多→肺癌高发。改造炉灶后发病率显著下降。
\end{tipbox}

\section{复习题}

\begin{enumerate}[leftmargin=*, itemsep=8pt]
    \item \textbf{【名词解释】}生物地球化学性疾病；碘缺乏病（IDD）
    \item \textbf{【名词解释】}地方性氟中毒；氟斑牙；氟骨症
    \item \textbf{【名词解释】}克山病；地方性砷中毒
    \item \textbf{【简答题】}简述生物地球化学性疾病的流行特征。
    \item \textbf{【简答题】}简述碘缺乏病的流行特征及预防措施。
    \item \textbf{【简答题】}简述地方性氟中毒的病区类型及预防措施。
    \item \textbf{【简答题】}简述氟的生理作用及双重健康效应。
    \item \textbf{【简答题】}简述慢性甲基汞中毒（水俣病）的病因及发病机制。
    \item \textbf{【简答题】}简述克山病的病因、病理特征及预防措施。
    \item \textbf{【简答题】}简述宣威肺癌高发的环境流行病学研究及其启示。
\end{enumerate}

\subsection*{参考答案要点}

\begin{enumerate}[leftmargin=*, itemsep=5pt]
    \item \textbf{生物地球化学性疾病}：因地球化学元素分布不均，摄入元素过多或过少引起的特异性疾病。\textbf{IDD}：从胚胎至成人期因碘摄入不足引起的一系列病症。

    \item \textbf{地方性氟中毒}：长期摄入过量氟引起的以氟骨症和氟斑牙为特征的慢性全身性疾病。

    \item \textbf{克山病}：与环境低硒有关的地方性心肌病。\textbf{地方性砷中毒}：长期摄入过量砷引起以皮肤损害为特征的慢性中毒。

    \item 两大特征：明显地区性分布；与环境中元素水平相关。影响因素：营养条件、生活习惯、多元素联合作用。

    \item 分布：山区$>$丘陵$>$平原$>$沿海，内陆$>$沿海。预防：碘盐、碘油、富碘食物。

    \item 三类病区：饮水型（改换水源+除氟）、燃煤污染型（改良炉灶）、饮砖茶型（低氟砖茶）。

    \item 适量：预防龋齿、促进骨骼钙化。过量：氟斑牙、氟骨症。缺乏：龋齿增多。呈U型剂量-效应关系。

    \item 来源：含汞废水→甲基化→食物链放大→人食鱼中毒。机制：甲基汞穿过血脑和胎盘屏障，损害中枢神经系统。

    \item 病因：低硒+柯萨奇病毒感染。病理：心肌变性坏死纤维化。预防：口服亚硒酸钠、改善膳食。

    \item 病因：室内燃煤产生PAHs。特点：女性发病率高于男性。启示：流行病学+干预研究整合范式。
\end{enumerate}

% ----- PPT参考题 -----
\subsection*{参考题（PPT课件补充——课程目标自测）}

\begin{enumerate}[leftmargin=*, itemsep=10pt]
    \item \textbf{【简答题】}什么是生物地球化学性疾病？其核心特征是什么？
    \item \textbf{【简答题】}生物地球化学性疾病有哪些主要流行特征？
    \item \textbf{【简答题】}我国主要防治的生物地球化学性疾病有哪些？（至少列出5种）
    \item \textbf{【简答题】}生物地球化学性疾病的流行受哪些影响因素的影响？举例说明。
    \item \textbf{【简答题】}简述生物地球化学性疾病的综合控制措施。
\end{enumerate}

\subsubsection*{参考答案}
\begin{enumerate}[leftmargin=*, itemsep=5pt]
    \item \textbf{定义}：由于地壳表面化学元素分布的不均匀性，使某些地区的水和/或土壤中某些元素过多或过少，当地居民通过饮水、食物等途径摄入这些元素过多或过少，而引起的特异性疾病。\textbf{核心特征}：呈明显地区性分布，与特定地理环境中某元素水平密切相关。
    \item 两大特征：①明显的地区性分布——山区、内陆等特定地理单元高发；②与环境中元素水平密切相关——摄入量过多或过少均可致病。
    \item 碘缺乏病（地方性甲状腺肿、地方性克汀病）、地方性氟中毒（饮水型/燃煤污染型/饮茶型）、地方性砷中毒（饮水型/燃煤污染型）、克山病（低硒相关地方性心肌病）、大骨节病（低硒相关骨关节病）。
    \item 三大影响因素：①营养条件——蛋白质充足可拮抗氟、砷毒性，维生素C促进氟排泄；②生活习惯——如燃煤烘烤食物、饮用砖茶等；③多种元素联合作用——如低碘+低硒协同加重碘缺乏病，高氟+低碘使氟中毒提前出现。
    \item 控制措施：①限制摄入——改换水源除氟/砷、改良炉灶、改变生活方式；②适量补充——食盐加碘（最安全有效）、口服亚硒酸钠补硒；③组织措施——建立防治机构、监测病情、动员群众参与。
\end{enumerate}

\newpage

% =====================================================================
%  CHAPTER 9-10: 公共场所卫生 & 环境质量评价
% =====================================================================
% =====================================================================
%  CHAPTER 9: 家用化学品卫生
% =====================================================================
\newpage